\subsection{Decomposition into irreducible factors in principal ideal domains}

We will now apply the results of VI, p.~18, relating to decomposition into irreducible elements, to principal ideal domains. By Prop. 2, an element \(p \neq 0\) of A is irreducible if and only if the ring \(\mathrm{A}/(p)\) is a field (I, p.~115, Cor. 1), that is if and only if the congruence \(ax \equiv b \pmod{p}\) admits a solution in A for each \(b \in \mathrm{A}\) and for each \(a \in \mathrm{A}\) which is not a multiple of \(p\) .

DEFINITION 2. --- Let A be an integral domain. A system of representatives of irreducible elements of A is a family \((p_{\alpha})\) of irreducible elements of A such that every irreducible element of A is an associate of precisely one \(p_{\alpha}\) .

THEOREM 2. --- Let A be a principal ideal domain and let \((p_{\alpha})\) be a system of representatives of irreducible elements of A. Then every nonzero element x of the field of fractions of A can be uniquely expressed in the form

\[
x = u \prod_ {\alpha} p _ {\alpha} ^ {n _ {\alpha}},\tag{1}
\]

where u is an invertible element of A, and where the \(n_{\alpha}\) are integers, all but finitely many of which are zero. For x to belong to A it is necessary and sufficient that all the \(n_{\alpha}\) be positive.

We will use the theorem about decomposition as a sum of irreducible elements (VI, p.~18, Th. 2), of which the above statement is only a translation. Since \(\mathcal{P}^*\) is a lattice ordered group it will be sufficient for us to show that every non empty set of principal ideals of A contains a maximal element, in order to check that the hypotheses of this theorem are indeed satisfied; now this follows from the next Lemma:

Lemma 1. --- Let A be a ring such that every left ideal of A is finitely generated. Then every nonempty set \(\Phi\) of left ideals of A, ordered by inclusion, has a maximal element.

By Zorn's Lemma (Set Theory, III, p.~154, Th. 2) it is enough to prove that \(\Phi\) is inductive. Now if \((\mathfrak{a}_{\lambda})\) is a totally ordered family of elements of \(\Phi\) then the union \(\mathfrak{a}\) of the ideals \(\mathfrak{a}_{\lambda}\) is a left ideal of \(\mathbf{A}\) , and so admits a finite system of generators \((a_i)_{1 \leqslant i \leqslant n}\) . Since each \(a_i\) belongs to an ideal \(\mathfrak{a}_{\lambda_i}\) , and since the family \((\mathfrak{a}_{\lambda})\) is totally ordered, the \(a_i\) all belong to the largest of the ideals \(\mathfrak{a}_{\lambda_i}\) , say \(\mathfrak{a}_{\mu}\) . Then \(\mathfrak{a} = \mathfrak{a}_{\mu}\) belongs to \(\Phi\) , which is thus indeed an inductive set.

Later we will study those rings B, called noetherian rings, such that every nonempty set of ideals of B contains a maximal element.

Remark. --- The family \((u, (n_{\alpha}))\) is called the decomposition of x into irreducible factors; by abuse of language, we also say that the formula (1) is the decomposition of x into irreducible factors. If \(x = u \prod_{\alpha} p_{\alpha}^{n_{\alpha}}\) and \(y = v \prod_{\alpha} p_{\alpha}^{m_{\alpha}}\) are

the decompositions of x and y into irreducible factors, then a necessary and sufficient condition for x to divide y is that \(n_{\alpha} \leqslant m_{\alpha}\) for all \(\alpha\) ; from this we deduce the formulae

\[
\operatorname * {g c d} (x, y) = \prod_ {\alpha} p _ {\alpha} ^ {\inf (n _ {\alpha}, m _ {\alpha})}\tag{2}
\]

\[
\operatorname{lcm} (x, y) = \prod_ {\alpha} p _ {\alpha} ^ {\sup \left(n _ {\alpha}, m _ {\alpha}\right)}.\tag{3}
\]

The property expressed by Th. 2 is true for a more general class of rings than principal ideal domains; we will study them later as unique factorisation domains; and we will see that polynomial rings and formal power series rings in arbitrarily many indeterminates are unique factorisation domains (Comm. Alg., VII, § 3).

